\chapter{Urutan Operasi Hitung}\label{ch:g7:priorities}

Berapa $3 + 4 \times 5$? Kalau kamu menghitung dari kiri ke kanan kamu
memperoleh $35$; kalau kamu mengalikan lebih dulu kamu memperoleh
$23$. Matematika memerlukan satu jawaban saja, jadi ada aturan lalu
lintas untuk perhitungan --- yaitu urutan pengerjaannya. Bab ini
menyatakan aturan itu lalu memperkenalkan kesamaan aljabar yang
pertama, yaitu sifat distributif.

\section{Aturan urutan pengerjaan}

\begin{theorem}[Aturan urutan pengerjaan]\label{thm:g7:priorities:rules}
Pada perhitungan tanpa tanda kurung:
\begin{enumerate}
  \item perkalian dan pembagian dikerjakan \emph{sebelum} penjumlahan
        dan pengurangan;
  \item operasi yang setingkat dikerjakan dari kiri ke kanan.
\end{enumerate}
Tanda kurung mengalahkan segalanya: yang ada di dalam kurung dihitung
lebih dulu, dimulai dari kurung yang paling dalam.
\end{theorem}

\begin{example}\label{ex:g7:priorities:basic}
Satu langkah per baris, sambil menggarisbawahi (dalam hati) apa yang dihitung:
\begin{align*}
3 + 4 \times 5 &= 3 + 20 = 23
&&\text{(perkalian lebih dulu)}\\
20 - 8 + 2 &= 12 + 2 = 14
&&\text{(setingkat: kiri ke kanan)}\\
18 \div 3 \times 2 &= 6 \times 2 = 12
&&\text{(setingkat: kiri ke kanan)}\\
(3 + 4) \times 5 &= 7 \times 5 = 35
&&\text{(kurung lebih dulu).}
\end{align*}
Perhatikan perangkapnya: $20 - 8 + 2$ itu \emph{bukan} $20 - 10$, dan
$18 \div 3 \times 2$ itu \emph{bukan} $18 \div 6$.
\end{example}

\begin{example}[Kurung bersarang]\label{ex:g7:priorities:nested}
\begin{align*}
50 - 2 \times \bigl(3 + (10 - 4)\bigr)
&= 50 - 2 \times (3 + 6) && \text{(kurung terdalam)}\\
&= 50 - 2 \times 9 && \text{(kurung yang tersisa)}\\
&= 50 - 18 && \text{(perkalian)}\\
&= 32 .
\end{align*}
\end{example}

\begin{remark}[Garis pecahan menyembunyikan tanda kurung]
Garis \omterm{def:g6:fractions:def}{pecahan} mengelompokkan \omterm{def:g4:fractions:def}{pembilang} dan \omterm{def:g4:fractions:def}{penyebutnya}, seolah-olah
masing-masing punya tanda kurung yang tidak terlihat:
\[
\frac{12 + 8}{4} = \frac{20}{4} = 5,
\qquad\text{bukan}\quad 12 + \frac{8}{4} = 14 .
\]
\end{remark}

\begin{method}[Menghitung dengan aman]\label{met:g7:priorities:safe}
\begin{enumerate}
  \item Kenalilah operasi yang harus dikerjakan lebih dulu (kurung
        terdalam, lalu $\times$ dan $\div$, lalu $+$ dan $-$);
  \item hitunglah \emph{hanya operasi itu}, lalu tulis ulang seluruh
        bentuknya dengan hasilnya di tempat itu;
  \item ulangi sampai tersisa satu bilangan --- satu operasi per
        baris, tanpa jalan pintas.
\end{enumerate}
\end{method}

\section{Sifat distributif}

\begin{theorem}[Sifat distributif]\label{thm:g7:priorities:distributivity}
Untuk semua bilangan $k$, $a$, $b$:
\[
k \times (a + b) = k \times a + k \times b,
\qquad
k \times (a - b) = k \times a - k \times b .
\]
\end{theorem}

\begin{proof}[Pembuktian dengan gambar]
Bilanglah luas persegi panjang selebar $k$ yang dibelah menjadi dua
persegi panjang yang lebih kecil dengan panjang $a$ dan $b$: sebagai
satu kesatuan luasnya $k \times (a+b)$; dalam dua bagian luasnya
$k \times a + k \times b$. Persegi panjang yang sama, luas yang sama.
\end{proof}

\begin{omfigure}
\begin{tikzpicture}[scale=0.75]
  \fill[omDef!14] (0,0) rectangle (4.2,1.8);
  \fill[omThm!14] (4.2,0) rectangle (6.4,1.8);
  \draw[thick] (0,0) rectangle (6.4,1.8);
  \draw[thick] (4.2,0) -- (4.2,1.8);
  \node[below, font=\small] at (2.1,0) {$a$};
  \node[below, font=\small] at (5.3,0) {$b$};
  \node[left, font=\small] at (0,0.9) {$k$};
  \node[font=\small] at (2.1,0.9) {$k \times a$};
  \node[font=\small] at (5.3,0.9) {$k \times b$};
\end{tikzpicture}

{\small Sifat distributif sebagai luas: persegi panjang yang besar
$k \times (a + b)$ adalah \omterm{def:g1:addition:def}{jumlah} kedua persegi panjang yang kecil.}
\end{omfigure}

\begin{example}[Kedua arahnya]\label{ex:g7:priorities:distributivity}
\emph{Menjabarkan} (dari kiri ke kanan) memudahkan hitungan di kepala:
\[
7 \times 103 = 7 \times (100 + 3) = 700 + 21 = 721,
\qquad
6 \times 98 = 6 \times (100 - 2) = 600 - 12 = 588 .
\]
\emph{Memfaktorkan} (dari kanan ke kiri) mengumpulkan faktor bersamanya:
\[
23 \times 12 + 23 \times 8 = 23 \times (12 + 8) = 23 \times 20 = 460 .
\]
\end{example}

\begin{example}[Menerjemahkan kata menjadi perhitungan]\label{ex:g7:priorities:words}
``\omterm{def:g1:addition:def}{Jumlah} dari $8$ dan \omterm{def:g2:mult:def}{hasil kali} $5$ dengan $3$'' adalah
$8 + 5 \times 3 = 23$. ``\omterm{def:g2:mult:def}{Hasil kali} $5$ dengan \omterm{def:g1:addition:def}{jumlah} dari $8$ dan
$3$'' adalah $5 \times (8 + 3) = 55$: kata ``\omterm{def:g2:mult:def}{hasil kali} \dots\ dengan
\omterm{def:g1:addition:def}{jumlah} dari\dots'' memasang tanda kurung di sekeliling \omterm{def:g1:addition:def}{jumlahnya}.
\end{example}

\section{Latihan}

\begin{exercise}[$\star$]\label{exo:g7:priorities:1}
Hitunglah, satu operasi per baris:
\[
5 + 3 \times 6, \qquad
(5 + 3) \times 6, \qquad
30 - 12 \div 4, \qquad
(30 - 12) \div 4 .
\]
\end{exercise}

\begin{exercise}[$\star$]\label{exo:g7:priorities:2}
Hitunglah:
\[
24 - 6 + 2, \qquad
24 \div 6 \times 2, \qquad
24 - 6 \times 2, \qquad
24 \div (6 \times 2).
\]
\end{exercise}

\begin{exercise}[$\star$]\label{exo:g7:priorities:3}
Hitunglah selangkah demi selangkah:
\[
7 + 3 \times (10 - 6), \qquad
40 - (2 + 3) \times 4, \qquad
60 \div \bigl(2 \times (7 - 4)\bigr).
\]
\end{exercise}

\begin{exercise}[$\star$]\label{exo:g7:priorities:4}
Hitunglah \omterm{def:g6:fractions:def}{pecahannya} (ada kurung yang tidak terlihat!):
\[
\frac{18 + 6}{8}, \qquad
\frac{30}{7 - 2}, \qquad
\frac{5 \times 8}{4 + 6} .
\]
\end{exercise}

\begin{exercise}[$\star$]\label{exo:g7:priorities:5}
Pakailah sifat distributif untuk menghitung di kepala:
\[
8 \times 102, \qquad
25 \times 99, \qquad
14 \times 12 + 14 \times 8, \qquad
37 \times 45 - 37 \times 35 .
\]
\end{exercise}

\begin{exercise}[$\star$]\label{exo:g7:priorities:6}
Tulislah perhitungannya dalam satu baris, dengan tanda kurung bila
perlu, lalu hitunglah: ``\omterm{ex:g1:subtraction:difference}{selisih} dari $50$ dan \omterm{def:g2:mult:def}{hasil kali} $6$ dengan
$7$''; ``\omterm{def:g2:mult:def}{hasil kali} dari \omterm{def:g1:addition:def}{jumlah} $4$ dan $5$ dengan \omterm{ex:g1:subtraction:difference}{selisih} $10$ dan
$8$''.
\end{exercise}

\begin{exercise}[$\star$]\label{exo:g7:priorities:7}
Amir membeli $3$ buku tulis seharga $2.50$ masing-masing dan $2$ pena
seharga $1.20$ masing-masing. Tulislah seluruh belanjanya dalam satu
bentuk, lalu hitunglah.
\end{exercise}

\begin{exercise}[$\star\star$]\label{exo:g7:priorities:8}
Pasanglah tanda kurung pada $4 + 2 \times 5 - 1$ agar hasilnya: (a)
$13$; (b) $29$; (c) $24$. (Salah satu dari ketiganya tidak memerlukan
tanda kurung sama sekali --- yang mana?)
\end{exercise}

\begin{exercise}[$\star\star$]\label{exo:g7:priorities:9}
Apakah $12 + 4 \times k$ dan $16 \times k$ sama untuk setiap bilangan
$k$? Ujilah dengan $k = 1$ dan $k = 2$, lalu simpulkan. Kesalahan
urutan pengerjaan yang mana yang membawa dari yang satu ke yang lain?
\end{exercise}

\begin{exercise}[$\star\star$]\label{exo:g7:priorities:10}
Karcis bioskop berharga $9$ euro, tetapi rombongan yang berisi $n$
orang membayar biaya tetap $12$ euro ditambah $6$ euro untuk tiap
orang.
\begin{enumerate}
  \item Tulislah bentuk untuk harga rombongan itu, lalu hitunglah
        untuk $n = 5$.
  \item Mulai dari berapa orang tarif rombongan lebih murah daripada
        $n$ karcis satuan? (Cobalah beberapa nilai $n$.)
\end{enumerate}
\end{exercise}

\begin{exercise}[$\star\star\star$]\label{exo:g7:priorities:11}
Dengan memakai masing-masing angka $1$, $2$, $3$, $4$ tepat sekali,
keempat operasi dan tanda kurung sesuka hati, tulislah bentuk yang
sama dengan $24$, dengan $10$, dan dengan $1$.
\end{exercise}

\section{Soal: Mengalikan seperti juru tulis Mesir}

\begin{problem}[{Soal akhir pekan --- trik hitung di kepala dari sifat
distributif, dan perkalian lipat dua--bagi dua yang dipakai selama
empat ribu tahun}]
\label{pb:g7:priorities:1}
Jauh sebelum ada kalkulator --- dan sebelum tabel perkalian dihafalkan
di sekolah --- para juru tulis Mesir dan petani Rusia mengalikan dua
bilangan mana pun hanya dengan \emph{melipatduakan, membagi dua dan
menjumlahkan}. Caranya tampak seperti sihir; sebenarnya itu sifat
distributif (\cref{thm:g7:priorities:distributivity}) yang bekerja
dengan tenaga penuh. Soal ini mengasah hitungan di kepalamu, lalu
mengajarkan cara kuno itu, lalu membuka kapnya untuk melihat dengan
tepat mengapa ia berhasil.

\medskip
\noindent\textbf{Bagian I --- Seni tidak menghitung.}
\begin{enumerate}
  \item Hitunglah di kepala, dengan menjabarkan seperti pada
        \cref{ex:g7:priorities:distributivity}:
        $7 \times 102$ dan $9 \times 98$.
  \item Hitunglah di kepala, dengan memfaktorkan:
        $17 \times 13 + 17 \times 7$ dan
        $45 \times 101 - 45$.
  \item Trik ``kali $5$'': untuk mengalikan dengan $5$, kalikan dengan
        $10$ lalu bagi dua. Hitunglah $5 \times 87$ dengan cara ini,
        lalu jelaskan mengapa triknya sah.
  \item Dua trik lagi untuk dijelaskan dan dipakai: mengalikan dengan
        $25$ dengan cara mengalikan dengan $100$ lalu membagi $4$
        (hitunglah $25 \times 32$); mengalikan dengan $99$ dengan cara
        mengalikan dengan $100$ lalu mengurangi sekali (hitunglah
        $99 \times 47$).
  \item Sebuah tembakan peringatan: sifat distributif menyebarkan
        sebuah faktor ke atas sebuah \emph{\omterm{def:g1:addition:def}{jumlah}}, tidak pernah ke
        atas \omterm{def:g2:mult:def}{hasil kali}. Bandingkan $8 \times (5 + 3)$ dengan
        $8 \times 5 + 8 \times 3$, lalu $8 \times (5 \times 3)$ dengan
        $(8 \times 5) \times (8 \times 3)$: pasangan mana yang cocok,
        dan pasangan mana yang tidak?
\end{enumerate}

\medskip
\noindent\textbf{Bagian II --- Cara lipat dua--bagi dua.}
Inilah resep untuk $13 \times 24$. Tulislah dua kolom. Kolom kiri:
mulailah dari $13$ lalu \emph{bagi dua}, dengan membuang \omterm{def:g3:division:remainder}{sisanya},
sampai $1$: \ $13,\ 6,\ 3,\ 1$. Kolom kanan: mulailah dari $24$ lalu
\emph{lipatduakan} sebanyak itu juga: \ $24,\ 48,\ 96,\ 192$. Coretlah
setiap baris yang bilangan \emph{kirinya} \omterm{def:g3:numbers:evenodd}{genap}, lalu jumlahkan
bilangan kanan yang selamat.
\begin{enumerate}[resume]
  \item Kerjakan resep itu untuk $13 \times 24$: baris mana yang
        dicoret, \omterm{def:g1:addition:def}{jumlah} berapa yang tersisa, dan apakah \omterm{def:g1:addition:def}{jumlah} itu
        sama dengan $13 \times 24$?
  \item Kalikan $21 \times 17$ dengan cara itu (bagi dua $21$,
        lipatduakan $17$), lalu periksa jawabannya dengan perkalian
        biasa.
  \item Kalikan $16 \times 25$ dengan cara itu. Berapa baris yang
        selamat, dan apa istimewanya kolom kiri yang bermula dari $1$,
        $2$, $4$, $8$, $16, \dots$?
  \item Sekarang \emph{mengapa} cara itu berhasil, selangkah demi
        selangkah. Berilah alasan setiap kesamaan dengan sifat
        distributif:
        \[
        13 \times 24 = (12 + 1) \times 24
        = 12 \times 24 + 24 = 6 \times 48 + 24 .
        \]
        Pesan moralnya: membagi dua yang kiri dan melipatduakan yang
        kanan mempertahankan \omterm{def:g2:mult:def}{hasil kalinya} --- tetapi ketika bilangan
        kirinya \emph{\omterm{def:g3:numbers:evenodd}{ganjil}}, satu salinan bilangan kanannya
        terpisah dan harus ditabung. Bilangan mana yang ditabung
        baris pertama tabel itu?
  \item Lanjutkan pembukuannya menuruni tabel:
        $6 \times 48 = 3 \times 96$, lalu
        $3 \times 96 = 1 \times 192 + 96$, lalu
        $1 \times 192 = 192$. Kumpulkan semua yang ditabung di
        sepanjang jalan, lalu jelaskan mengapa bilangan yang ditabung
        itu tepat bilangan kolom kanan pada baris yang bilangan
        kirinya \omterm{def:g3:numbers:evenodd}{ganjil} --- yaitu baris yang disimpan resepnya.
\end{enumerate}

\medskip
\noindent\textbf{Bagian III --- Membuka kapnya: \omterm{def:g1:addition:def}{jumlah} dari
lipatan dua.}
\begin{enumerate}[resume]
  \item Pada pertanyaan 6, baris yang selamat bernilai $1$, $4$
        dan $8$ kali lipatan dua dari $24$: caranya diam-diam menulis
        $13 \times 24 = (1 + 4 + 8) \times 24$. Periksalah bahwa
        $1 + 4 + 8 = 13$, lalu carilah lipatan dua mana ($1$, $2$,
        $4$, $8$, $16$) yang disimpan caranya untuk $21$ pada
        pertanyaan 7.
  \item Tulislah $25$ dan $42$ sebagai \omterm{def:g1:addition:def}{jumlah} bilangan dari daftar
        lipatan dua $1, 2, 4, 8, 16, 32$, dengan memakai masing-masing
        paling banyak sekali.
  \item Jelaskan mengapa \emph{setiap} bilangan cacah dapat ditulis
        dengan cara ini: ambillah bilangan lipatan dua terbesar yang
        muat, kurangkan, lalu ulangi --- mengapa ini pasti berhenti?
        Tulislah $100$ sebagai \omterm{def:g1:addition:def}{jumlah} semacam itu.
  \item Para juru tulis Mesir melakukan persis ini, empat ribu tahun
        yang lalu: untuk mengalikan $19 \times 35$, mereka membuat
        tabel lipatan dua dari $35$ ($1 \to 35$, $2 \to 70$,
        $4 \to 140$, $8 \to 280$, $16 \to 560$), memilih baris yang
        berjumlah $19$, lalu menjumlahkannya. Kerjakan perhitungannya
        dengan cara mereka.
  \item Penutupnya: komputer modern, yang membilang dengan angka $0$
        dan $1$, mengalikan pada dasarnya seperti para juru tulis itu
        --- dengan lipatan dua dan penjumlahan. Tiga penjumlahan yang
        mana yang akan dilakukan komputer untuk $13 \times 24$?
        (Mengapa setiap bilangan terpecah menjadi lipatan dua, dan apa
        hubungan angka $0$ dan $1$ dengan hal itu, adalah soal akhir
        pekan bab Perpangkatan: \cref{pb:g8:powers:1}.)

\end{enumerate}
\end{problem}
